\documentclass[12pt]{letter}
\usepackage[utf8]{inputenc}
\usepackage[T1]{fontenc}
\usepackage[margin=1in]{geometry}
\usepackage{siunitx}
\usepackage{hyperref}

\sisetup{per-mode=symbol}

\signature{Jonathan Mace\\
Microsoft Research, Redmond, WA, United States}

\address{Jonathan Mace\\
Microsoft Research\\
Redmond, WA\\
United States}

\date{\today}

\begin{document}

\begin{letter}{%
  Editor\\
  Journal of the Acoustical Society of America\\
  Acoustical Society of America%
}

\opening{Dear Editor,}

We are pleased to submit the enclosed manuscript, \emph{``A Constrained-Bubble
Model for Gas-Pocket Mechanotransduction in the Human Abdomen: Theoretical
Framework and Sensitivity Analysis,''} for consideration as an original
research article in the \emph{Journal of the Acoustical Society of America}.

This paper investigates a previously unexplored acoustic mechanism:
intraluminal gas pockets in the gastrointestinal tract acting as
impedance-matching transducers that convert airborne infrasonic pressure
oscillations into local intestinal wall strain.  We develop a
constrained-bubble model---Minnaert dynamics modified for the elastic
intestinal wall and cylindrical lumen geometry---and show that gas pockets of
\SIrange{5}{100}{\milli\litre} amplify tissue displacement as compliant
inclusions within an otherwise nearly incompressible medium.

The principal findings are:
\begin{itemize}
  \item Gas pockets provide a 35--100$\times$ efficiency advantage over
    airborne coupling to whole-cavity flexural modes, representing the most
    plausible acoustic pathway identified to date for infrasound-induced GI
    mechanotransduction.
  \item The SPL threshold for activating \textsc{Piezo} mechanosensitive ion
    channels ranges from \SI{111}{\decibel} (\SI{100}{\milli\litre} pocket)
    to \SI{120}{\decibel} (\SI{5}{\milli\litre}), within the range of
    occupational and environmental exposures.
  \item Under the idealised sealed-pocket assumption, a Monte Carlo
    population model ($N = 10\,000$) places most sampled individuals
    above the illustrative \textsc{Piezo} threshold at
    \SI{120}{\decibel}, highlighting gas volume and acoustic sealing as
    the dominant determinants of susceptibility.
\end{itemize}
This work is a companion to a manuscript currently under review at the
\emph{Journal of Sound and Vibration}, which establishes the inadequacy of
whole-cavity airborne coupling (the ``brown note'' hypothesis).  The present
paper is fully self-contained and addresses a distinct physical mechanism
--- localised gas-pocket transduction rather than global cavity resonance.

We believe the novelty of this contribution is significant: to our knowledge,
this is the first analysis of constrained-bubble dynamics in bowel gas under
infrasonic excitation.  We have been careful to address the acoustic
short-circuit caveat honestly throughout the manuscript, noting that the
surrounding tissue compliance limits the effective impedance contrast at low
frequencies and that our estimates therefore represent upper bounds on the
gas-pocket pathway's efficiency.

The manuscript has not been published or submitted elsewhere.  All authors
have approved the manuscript and agree with its submission to JASA.

We suggest the following subject classifications: Biomedical Acoustics;
Nonlinear Acoustics; Underwater and Bubble Acoustics.

Thank you for considering this manuscript.  We look forward to your response.

\closing{Yours sincerely,}

\end{letter}
\end{document}
